\hypertarget{appendix-c83016554b1e}{}
\section*{Confidence intervals and sequential evidence}
\section{Confidence intervals and sequential evidence}\label{confidence-intervals-and-sequential-evidence}

\textbf{Curriculum depth:} Developed · \textbf{First seen:} Phase 10 · \textbf{Latest development:} Phase 10

\subsection{Definition and mechanism}\label{definition-and-mechanism}

A confidence interval is a procedure with repeated-sample coverage under stated assumptions; it is not a posterior probability assigned to a fixed parameter. A normal interval for a mean uses a known population standard deviation or an approximation justified by the central limit theorem. With unknown scale and IID normal observations, the Student t distribution accounts for estimating the standard deviation. Heavy tails, sample dependence, multiple comparisons, and optional stopping can invalidate simple textbook guarantees. Adaptive algorithms that stop when a fixed-sample interval looks narrow need sequentially valid methods or a separately justified stopping rule. Report the target, sampling unit, dependence structure, estimator bias, interval procedure, stopping policy, and independent-run count. A narrow interval around a biased estimator measures precision about its expectation, not correctness of the target.

\subsection{Prerequisites}\label{prerequisites}

\begin{itemize}
\tightlist
\item
  \hyperref[appendix-a2360b435383]{Common probability distributions and their assumptions}
\item
  \hyperref[appendix-54c383df9992]{Expectation, covariance, and dependence}
\item
  \href{}{Monte Carlo rendering estimators and variance (04-Visual-Computing-II, appendix-031e02bf28af)}
\end{itemize}

\subsection{Used by}\label{used-by}

\begin{itemize}
\item
  This node is a current endpoint in the supplied Phase 9--10 path.
\item
  \hyperref[appendix-b9db93e28371]{Reproducible numerical benchmarking} \#\# Full teaching explanation
\item
  \href{Continuous\%20Notes/Computing\%20Foundations\%20-\%20From\%20Bits\%20to\%20C++\%20and\%20Systems\%20APIs\#10.\%20Inference\%20and\%20error\%20metrics\%20need\%20explicit\%20assumptions}{Phase 10 teaching section}
\item
  \href{}{Visual Recall/10 - Probability and Monte Carlo - Visual Recall.md (01-Foundations-I, appendix-cb5af353ef77)} · visual recall
\end{itemize}

\subsection{Source coverage}\label{source-coverage}

Supplied lesson parts: 960 (T06 · M058), 961 (T06 · M058).
